
\documentclass{article}
\usepackage{colortbl}
\usepackage{makecell}
\usepackage{multirow}
\usepackage{supertabular}

\begin{document}

\newcounter{utterance}

\centering \large Interaction Transcript for game `cladder', experiment `full\_v1.5\_default', episode 9763 with Qwen3.5{-}4B.
\vspace{24pt}

{ \footnotesize  \setcounter{utterance}{1}
\setlength{\tabcolsep}{0pt}
\begin{supertabular}{c@{$\;$}|p{.15\linewidth}@{}p{.15\linewidth}p{.15\linewidth}p{.15\linewidth}p{.15\linewidth}p{.15\linewidth}}
    \# & \multicolumn{2}{c}{Player} && \multicolumn{2}{c}{Game Master} \\
    \hline

    \theutterance \stepcounter{utterance}  
    & & & \multicolumn{4}{p{0.6\linewidth}}{
        \cellcolor[rgb]{0.9,0.9,0.9}{
            \makecell[{{p{\linewidth}}}]{
                \texttt{\tiny{[P1$\langle$GM]}}
                \texttt{You are an expert in causal inference. The following question is not a typical commonsense query, but rather a meticulously designed question created by a professor specializing in causal inference, intended to assess the students' mastery of the course content.} \\
\texttt{Imagine a self{-}contained, hypothetical world with only the following conditions, and without any unmentioned factors or causal relationships: Yupt has a direct effect on jyka and kwox. Jyka has a direct effect on kwox. We know that yupt causes jyka. yupt and jyka causes kwox. We observed an individual is yupt. Would an individual is kwox if not jyka instead of jyka?} \\
\texttt{Start your answer with 'Yes' or 'No', followed by additional reasoning or evidence to support your explanation.} \\
            }
        }
    }
     \\ \\

    \theutterance \stepcounter{utterance}  
    & \multicolumn{4}{p{0.6\linewidth}}{
        \cellcolor[rgb]{0.561,0.737,0.561}{
            \makecell[{{p{\linewidth}}}]{
                \texttt{\tiny{[P1$\rangle$GM]}}
                \texttt{Yes} \\
\\ 
\texttt{**Reasoning:**} \\
\\ 
\texttt{To determine the answer, we must analyze the causal structure defined in the hypothetical world and apply the rules of causal inference regarding counterfactuals.} \\
\\ 
\texttt{1.  **Causal Structure**:} \\
\texttt{    *   \$Yupt \textbackslash{}rightarrow Jyka\$} \\
\texttt{    *   \$Yupt \textbackslash{}rightarrow Kwox\$} \\
\texttt{    *   \$Jyka \textbackslash{}rightarrow Kwox\$} \\
\texttt{    *   This forms a "fork" or "collider" structure depending on the direction of interest, but specifically, \$Kwox\$ is a common effect of both \$Yupt\$ and \$Jyka\$.} \\
\\ 
\texttt{2.  **The Observation**:} \\
\texttt{    *   We observed an individual where \$Yupt = 1\$ (True).} \\
\\ 
\texttt{3.  **The Counterfactual Question**:} \\
\texttt{    *   The question asks: "Would an individual be \$Kwox\$ if not \$Jyka\$ instead of \$Jyka\$?"} \\
\texttt{    *   Formally, this asks for the counterfactual outcome: \$Kwox\_\{\textbackslash{}text\{not \} Jyka\} \textbackslash{}mid Yupt=1\$.} \\
\texttt{    *   This is asking: If we intervene to set \$Jyka = 0\$ (block the effect of \$Jyka\$ on \$Kwox\$), while keeping \$Yupt = 1\$, does \$Kwox\$ still occur?} \\
\\ 
\texttt{4.  **Causal Mechanism Analysis**:} \\
\texttt{    *   Since \$Yupt\$ has a **direct effect** on \$Kwox\$, the occurrence of \$Kwox\$ is guaranteed by the presence of \$Yupt\$, regardless of the status of \$Jyka\$.} \\
\texttt{    *   The path \$Yupt \textbackslash{}rightarrow Kwox\$ remains active even if the path \$Yupt \textbackslash{}rightarrow Jyka \textbackslash{}rightarrow Kwox\$ is blocked by setting \$Jyka = 0\$.} \\
\texttt{    *   In causal diagrams, \$Kwox\$ is a collider on the path \$Jyka \textbackslash{}rightarrow Kwox \textbackslash{}leftarrow Yupt\$. However, the question is not about "explaining away" (which happens when we observe \$Kwox\$ and ask about \$Yupt\$ and \$Jyka\$); it is about an **intervention** on \$Jyka\$.} \\
\texttt{    *   When we intervene (\$do(Jyka=0)\$), we sever the influence of \$Jyka\$ on \$Kwox\$. However, the direct causal arrow from \$Yupt\$ to \$Kwox\$ is not severed.} \\
\\ 
\texttt{5.  **Conclusion**:} \\
\texttt{    *   Because \$Yupt\$ directly causes \$Kwox\$, the individual will still be \$Kwox\$ even if \$Jyka\$ is set to false. The direct causal link ensures the outcome persists.} \\
\\ 
\texttt{Therefore, the answer is **Yes**.} \\
            }
        }
    }
    & & \\ \\

    \theutterance \stepcounter{utterance}  
    & & & \multicolumn{2}{p{0.3\linewidth}}{
        \cellcolor[rgb]{0.9,0.9,0.9}{
            \makecell[{{p{\linewidth}}}]{
                \texttt{\tiny{[GM$|$GM]}}
                \texttt{Yes} \\
            }
        }
    }
    & & \\ \\

    \theutterance \stepcounter{utterance}  
    & & & \multicolumn{2}{p{0.3\linewidth}}{
        \cellcolor[rgb]{0.9,0.9,0.9}{
            \makecell[{{p{\linewidth}}}]{
                \texttt{\tiny{[GM$|$GM]}}
                \texttt{no} \\
            }
        }
    }
    & & \\ \\

    \theutterance \stepcounter{utterance}  
    & & & \multicolumn{2}{p{0.3\linewidth}}{
        \cellcolor[rgb]{0.9,0.9,0.9}{
            \makecell[{{p{\linewidth}}}]{
                \texttt{\tiny{[GM$|$GM]}}
                \texttt{game\_result = LOSE} \\
            }
        }
    }
    & & \\ \\

\end{supertabular}
}

\end{document}
